\section{Limitations}

The results constrain zero multiplicities but do not prove the Riemann hypothesis. Sharpness in Theorem~\ref{thm:sharp} is sharpness of a scalar relaxation: the extremal distributions need not arise from the analytic evaluation vectors of the zeta function. Further improvement from the same two-trace input would require additional geometric or arithmetic information.

\section*{Acknowledgments and disclosure}
The analytic input and its formal development are due to the author of \cite{Claude2026} and the teams acknowledged there. Zach Waddle directed the present project. OpenAI's GPT-5.6 Pro assisted with derivation, checking, source inspection, exact-certificate generation, and manuscript preparation. Independent specialist review and a compiled Lean extension remain outstanding.

\begin{thebibliography}{99}
\small
\bibitem{Claude2026} Claude, \emph{More than two thirds of the zeros of the Riemann zeta function lie on the critical line}, preprint, 2026; accompanying Lean repository \texttt{anthropics/zeta-23-lean}, commit \path{3635e74826a4c1fcece7d1cd2b6fa75e43a00510}.
\bibitem{GDL2025} F. Gon\c{c}alves, D. de Laat, and N. Leijenhorst, \emph{Multiplicity of nontrivial zeros of primitive $L$-functions via higher-level correlations}, Math. Comp. 94 (2025), 2041--2058.
\bibitem{STTB2022} A. Simoni\v{c}, T. Trudgian, and C. L. Turnage-Butterbaugh, \emph{Some explicit and unconditional results on gaps between zeroes of the Riemann zeta-function}, Trans. Amer. Math. Soc. 375 (2022), 3239--3265.
\bibitem{Montgomery1973} H. L. Montgomery, \emph{The pair correlation of zeros of the zeta function}, Proc. Sympos. Pure Math. 24 (1973), 181--193.
\bibitem{CGG1998} J. B. Conrey, A. Ghosh, and S. M. Gonek, \emph{Simple zeros of the Riemann zeta-function}, Proc. London Math. Soc. 76 (1998), 497--522.
\end{thebibliography}
